\section{Подробный статистический анализ линейной модели}

\subsection{Оценки коэффициентов}

Для линейной модели \(y = a + bx\) получены следующие оценки:

\[
\hat{a} = -24.358, \quad \hat{b} = 0.9879.
\]

\subsection{Прогнозные значения и остатки}

Прогнозные значения \(\hat{y}_i\) и остатки \(e_i = y_i - \hat{y}_i\) для первых десяти наблюдений приведены в таблице~\ref{tab:residuals}.

\begin{table}[H]
\centering
\caption{Прогнозные значения и остатки линейной модели (первые 10 строк)}
\label{tab:residuals}
\begin{tabular}{cccc}
\toprule
\(i\) & \(y_i\) & \(\hat{y}_i\) & \(e_i\) \\
\midrule
1 & 361.000 & 347.099 & 13.901 \\
2 & 419.000 & 378.712 & 40.288 \\
3 & 385.000 & 412.301 & \(-27.301\) \\
4 & 381.000 & 413.289 & \(-32.289\) \\
5 & 389.000 & 415.265 & \(-26.265\) \\
6 & 390.000 & 418.229 & \(-28.229\) \\
7 & 418.000 & 426.132 & \(-8.132\) \\
8 & 420.000 & 429.096 & \(-9.096\) \\
9 & 428.000 & 433.047 & \(-5.047\) \\
10 & 407.000 & 435.023 & \(-28.023\) \\
\bottomrule
\end{tabular}
\end{table}

Сумма остатков \(\sum e_i \approx 0\) (с точностью до вычислительной погрешности), что является свойством МНК-оценок.

\subsection{Остаточная сумма квадратов и оценка дисперсии ошибки}

Остаточная сумма квадратов:

\[
RSS = \sum_{i=1}^{n} e_i^2 = 8\,765\,550.312.
\]

Несмещённая оценка дисперсии ошибки:

\[
s^2 = \frac{RSS}{n-2} = \frac{8\,765\,550.312}{52} = 168\,568.275,
\]
\[
s = \sqrt{s^2} = 410.571.
\]

Значение \(s = 410.57\) характеризует среднюю величину разброса наблюдений относительно линии регрессии. По сравнению со средним значением \(y\) (\(\bar{y} = 1090.96\)) это составляет около 37.6\%, что указывает на умеренную точность модели.

\subsection{Доверительные интервалы для коэффициентов}

Стандартные ошибки коэффициентов:

\[
SE(\hat{b}) = \frac{s}{\sqrt{S_{xx}}} = \frac{410.571}{\sqrt{26\,186\,057.926}} = 0.0802,
\]
\[
SE(\hat{a}) = s\sqrt{\frac{1}{n} + \frac{\bar{x}^2}{S_{xx}}} = 410.571 \cdot \sqrt{\frac{1}{54} + \frac{1128.963^2}{26\,186\,057.926}} = 106.425.
\]

Квантиль распределения Стьюдента для уровня значимости \(\alpha = 0.05\) и \(\nu = n - 2 = 52\) степеней свободы:

\[
t_{0.975;\,52} = 2.0066.
\]

95\%-е доверительные интервалы:

\[
\hat{a} \pm t_{0.975;\,52} \cdot SE(\hat{a}) = -24.358 \pm 2.0066 \cdot 106.425 = [-237.916;\; 189.201],
\]
\[
\hat{b} \pm t_{0.975;\,52} \cdot SE(\hat{b}) = 0.9879 \pm 2.0066 \cdot 0.0802 = [0.827;\; 1.149].
\]

Доверительный интервал для коэффициента \(\hat{a}\) содержит ноль, что говорит о том, что свободный член линейной регрессии \textbf{статистически незначим} на уровне \(\alpha = 0.05\). Это ожидаемо, поскольку значение \(x = 0\) лежит далеко за пределами наблюдаемого диапазона.

Доверительный интервал для коэффициента \(\hat{b}\) \textbf{не содержит ноль}, что свидетельствует о статистической значимости линейной зависимости между \(x\) и \(y\).

\subsection{Проверка гипотезы о значимости коэффициента при \(x\)}

Формальная проверка гипотезы:

\[
H_0: b = 0 \quad \text{(линейная зависимость отсутствует)},
\]
\[
H_1: b \neq 0 \quad \text{(линейная зависимость присутствует)}.
\]

Уровень значимости: \(\alpha = 0.05\).

Наблюдаемое значение статистики Стьюдента:

\[
t_{\text{набл}} = \frac{\hat{b}}{SE(\hat{b})} = \frac{0.9879}{0.0802} = 12.313.
\]

Критическое значение: \(t_{\text{крит}} = t_{0.975;\,52} = 2.0066\).

Так как \(|t_{\text{набл}}| = 12.313 > 2.0066 = t_{\text{крит}}\), нулевая гипотеза \(H_0\) \textbf{отвергается} на уровне значимости \(\alpha = 0.05\).

Эквивалентно, с использованием p-value:

\[
p\text{-value} = 2 \cdot P(T_{52} \geq 12.313) < 0.0001.
\]

Поскольку \(p\text{-value} < \alpha = 0.05\), гипотеза \(H_0\) отвергается.

\subsection{Вывод о значимости}

Коэффициент при \(x\) в линейной модели является \textbf{статистически значимым} на уровне \(\alpha = 0.05\). Это подтверждает наличие линейной зависимости между \(x\) и \(y\): с ростом фактора \(x\) результат \(y\) в среднем возрастает. Данный вывод согласуется с тем, что доверительный интервал для \(\hat{b}\) не содержит ноль.
